% ===================================================================
%  ਸਾਰਣੀ ਨੰ. 12 — ਸਟੇਸ਼ਨ। — ਰੇਲ ਗੱਡੀਆਂ। — ਜਹਾਜ਼।
%  Traduction 2026 d'apres texto/io, controlee sur texto/fr, texto/en
%  et texto/hi. Consigne : voir texto/pa/10-sarni-01.tex.
% ===================================================================

%%K t12-apar-1 apar
\VUsaut{4.0mm}
\VUtitre{15.0pt}{60}{ਸਾਰਣੀ ਨੰ. 12}
\VUsaut{2.4mm}
\VUpk{11.4pt}{40}{ਸਟੇਸ਼ਨ। --- ਰੇਲ ਗੱਡੀਆਂ। --- ਜਹਾਜ਼।}
\VUsaut{3.4mm}
\textsc{ਟਿੱਪਣੀ।} --- \textit{ਹਰ ਦੇਸ ਵਿਚ ਵੇਲੇ ਦੀ ਗਿਣਤੀ ਵੱਖਰੀ ਹੈ, ਇਸੇ ਲਈ ਅਸੀਂ ਮਿਸਾਲ
ਲਈ} \VUgras{ਫ਼ਰਾਂਸੀਸੀ ਸੂਚੀ} \textit{ਦੇ ਵੇਲੇ ਲਏ ਹਨ।}

%%K t12-01-1 p
1. --- ਮੈਂ ਹੁਣੇ ਜਰਮਨੀ ਦਾ ਇਕ ਸਫ਼ਰ ਕੀਤਾ ਹੈ। ਨਿਕਲਣ ਤੋਂ ਪਹਿਲਾਂ ਮੈਂ ਗੱਡੀਆਂ ਦੇ ਛੁੱਟਣ
ਦਾ ਵੇਲਾ ਜਾਣਨ ਲਈ ਇਕ \VUgras{ਵੇਲੇ ਦੀ ਸੂਚੀ} \textsuperscript{(15)} ਖ਼ਰੀਦੀ ਸੀ। ਇਹ ਵੇਖ
ਕੇ ਕਿ ਸਭ ਤੋਂ ਸੌਖੀ ਗੱਡੀ ਪੈਰਿਸ ਤੋਂ ਰਾਤ ਨੌਂ ਵਜੇ ਛੁੱਟਦੀ ਹੈ ਤੇ ਸਟ੍ਰਾਸਬੂਰਗ ਇਕ ਵੱਜ ਕੇ
ਤੇਰਾਂ ਮਿੰਟ ਉੱਤੇ ਪਹੁੰਚਦੀ ਹੈ, ਮੈਂ ਆਪਣੇ ਸਫ਼ਰ ਦੀ ਤਿਆਰੀ ਕੀਤੀ, ਤੇ ਅਗਲੇ ਦਿਨ ਸਟੇਸ਼ਨ ਤਕ
ਗੱਡੀ ਕਰਾ ਲਈ।

%%K t12-02-1 p
2. --- ਜਦੋਂ ਮੈਂ ਤੁਰਨ ਦੇ ਵੱਡੇ ਕਮਰੇ ਵਿਚ ਵੜਿਆ, ਇਕ ਬੁੱਢੇ ਪਿੰਡ ਦੇ \VUgras{ਪਾਦਰੀ}
\textsuperscript{(2)} ਦੇ ਪਿੱਛੇ, ਜੋ ਹੱਥ ਵਿਚ ਆਪਣਾ \VUgras{ਸਫ਼ਰੀ ਥੈਲਾ}
\textsuperscript{(3)} ਚੁੱਕੀ ਸੀ, ਤਦੋਂ ਤਕ ਬਹੁਤ ਸਾਰੇ \VUgras{ਮੁਸਾਫ਼ਰ}
\textsuperscript{(1)} ਆ ਚੁੱਕੇ ਸਨ।

%%K t12-02-2 p
ਮੈਨੂੰ ਇਕ \VUgras{ਤਾਰ} \textsuperscript{(5)} ਘੱਲਣੀ ਸੀ, ਇਸੇ ਲਈ ਮੈਂ ਇਕ
\VUgras{ਨਿਗਰਾਨ} \textsuperscript{(6)} ਕੋਲੋਂ \VUgras{ਤਾਰ-ਘਰ} \textsuperscript{(4)}
ਪੁੱਛ ਲਿਆ।

%%K t12-03-1 p
3. --- \VUgras{ਉਡੀਕ-ਕਮਰੇ} \textsuperscript{(7)} ਵਿਚ ਜਾਣ ਤੋਂ ਪਹਿਲਾਂ ਮੈਂ ਆਉਣ-ਜਾਣ ਦਾ
ਇਕ \VUgras{ਟਿਕਟ} \textsuperscript{(9)} ਲੈਣ ਗਿਆ।

%%K t12-04-1 p
4. --- \VUgras{ਟਿਕਟ-ਖਿੜਕੀ} \textsuperscript{(8)} ਉੱਤੇ ਮੈਨੂੰ ਦੇਰ ਨਾ ਲੱਗੀ, ਕਿਉਂਕਿ
ਉੱਥੇ ਭੀੜ ਘੱਟ ਸੀ। ਛੁੱਟੀ ਉੱਤੇ ਜਾਂਦੇ ਇਕ \VUgras{ਸਿਪਾਹੀ} \textsuperscript{(10)} ਨੇ
ਮਿਹਰਬਾਨੀ ਕਰ ਕੇ ਮੈਨੂੰ ਆਪਣੀ ਵਾਰੀ ਦੇ ਦਿੱਤੀ, ਸੋ ਮੈਨੂੰ ਇੰਨਾ ਵੇਲਾ ਮਿਲ ਗਿਆ ਕਿ
\VUgras{ਸਟੇਸ਼ਨ ਮਾਸਟਰ} \textsuperscript{(13)} ਕੋਲੋਂ ਕੁਝ ਗੱਲਾਂ ਪੁੱਛ ਲਵਾਂ, ਜੋ
\VUgras{ਪੁਲਿਸ ਅਫ਼ਸਰ} \textsuperscript{(12)} ਤੇ ਡਿਊਟੀ ਉੱਤੇ ਖੜੇ \VUgras{ਜ਼ਾਂਦਾਰਮ}
\textsuperscript{(11)} ਨਾਲ ਗੱਲਾਂ ਕਰ ਰਹੇ ਸਨ।

%%K t12-tit-1 sub
\VUpk{10.2pt}{40}{ਸਾਮਾਨ ਦੀ ਬੁਕਿੰਗ।}
\VUsaut{3.0mm}

%%K t12-05-1 p
5. --- \VUgras{ਕਿਤਾਬਾਂ ਦੇ ਖੋਖੇ} \textsuperscript{(14)} ਤੋਂ ਇਕ ਅਖ਼ਬਾਰ ਖ਼ਰੀਦ ਕੇ ਮੈਂ
ਆਪਣਾ \VUgras{ਸਾਮਾਨ} \textsuperscript{(17)} ਤੁਲਵਾਇਆ — ਇਕ \VUgras{ਕੁਲੀ}
\textsuperscript{(16)} ਉਸ ਨੂੰ ਹੁਣੇ ਇਕ \VUgras{ਸਾਮਾਨ ਦੀ ਰੇੜ੍ਹੀ}
\textsuperscript{(19)} ਉੱਤੇ ਲਿਆਇਆ ਸੀ। \VUgras{ਬਾਬੂ} \textsuperscript{(22)}, ਜੋ
\VUgras{ਤੋਲ-ਜੰਤਰ} \textsuperscript{(21)} ਦੇ ਇੰਚਾਰਜ ਹਨ, ਉਨ੍ਹਾਂ ਮੈਨੂੰ ਦੱਸਿਆ ਕਿ ਮੇਰਾ
ਸੰਦੂਕ ਪੈਂਤੀ ਕਿੱਲੋ ਦਾ ਹੈ, ਜਿਸ ਨਾਲ ਪੰਜ ਕਿੱਲੋ ਦਾ \VUgras{ਵਾਧਾ} ਬਣਦਾ ਹੈ, ਜਿਸ ਲਈ
ਮੈਨੂੰ \VUgras{ਸਾਮਾਨ ਦੀ ਖਿੜਕੀ} \textsuperscript{(23)} ਉੱਤੇ ਵਾਧੂ ਮਹਿਸੂਲ ਦੇਣਾ ਪਵੇਗਾ।

%%K t12-06-1 p
6. --- ਕਿਉਂਕਿ ਮੈਂ ਵੇਲੇ ਤੋਂ ਪਹਿਲਾਂ ਸੀ, ਮੈਨੂੰ ਸਟੇਸ਼ਨ ਵਿਚ ਮੁਸਾਫ਼ਰਾਂ ਦੀ ਚਹਿਲ-ਪਹਿਲ ਵੇਖਣ
ਦਾ ਵੇਲਾ ਮਿਲ ਗਿਆ। ਕੁਝ \VUgras{ਅਮਾਨਤ-ਘਰ} \textsuperscript{(25)} ਵਿਚ ਨਿੱਕੇ ਥੈਲੇ ਰੱਖ
ਜਾਂ ਲੈ ਰਹੇ ਸਨ; ਦੂਜੇ \VUgras{ਵੇਲੇ ਦੀਆਂ ਸੂਚੀਆਂ} \textsuperscript{(26)} ਵੇਖ ਰਹੇ ਸਨ।
ਆਉਂਦੇ ਮੁਸਾਫ਼ਰ \VUgras{ਨਿਕਾਸ} \textsuperscript{(32)} ਵੱਲ ਕਾਹਲੀ ਕਰਦੇ ਸਨ, ਹੱਥ ਵਿਚ ਆਪਣੇ
\VUgras{ਅਟੈਚੀ} \textsuperscript{(24)} ਚੁੱਕੀ। ਕੁਝ \VUgras{ਸਾਮਾਨ ਦੀ ਸਪੁਰਦਗੀ}
\textsuperscript{(27)} ਉੱਤੇ ਉਡੀਕ ਕਰਦੇ ਤੇ ਆਪਣੀ \VUgras{ਰਸੀਦ} \textsuperscript{(29)}
ਵਿਖਾਉਂਦੇ, ਆਪਣੇ ਸੰਦੂਕ, \VUgras{ਪੇਟੀਆਂ} \textsuperscript{(28)} ਜਾਂ
\VUgras{ਟੋਕਰੀਆਂ} \textsuperscript{(30)} ਲੈਣ ਲਈ।

%%K t12-07-1 p
7. --- \VUgras{ਚੁੰਗੀ ਦੇ ਕਾਰਿੰਦੇ} \textsuperscript{(33)} ਧਿਆਨ ਨਾਲ ਗੱਠੜੀਆਂ ਪਰਖਦੇ ਸਨ
ਕਿ ਕਿਤੇ ਕੁਝ ਦੱਸਣ ਜੋਗ ਤਾਂ ਨਹੀਂ।

%%K t12-08-1 p
8. --- ਕੁਝ ਮੁਸਾਫ਼ਰ \VUgras{ਖਾਣ-ਪੀਣ ਦੇ ਕਮਰੇ} \textsuperscript{(34)} ਵੱਲ ਗਏ, ਜਿੱਥੇ
ਇਕ \VUgras{ਬੈਰਾ} \textsuperscript{(35)} ਉਨ੍ਹਾਂ ਨੂੰ ਫੁਰਤੀ ਨਾਲ ਵਰਤਾ ਰਿਹਾ ਸੀ। ਉਨ੍ਹਾਂ
ਨੂੰ ਕਾਹਲੀ ਕਰਨੀ ਸੀ, ਕਿਉਂਕਿ \VUgras{ਗੱਡੀ} \textsuperscript{(36)}, ਜੋ ਐਕਸਪ੍ਰੈਸ ਸੀ,
ਸਟੇਸ਼ਨ ਉੱਤੇ ਵੱਧ ਦੇਰ ਨਹੀਂ ਰੁਕਦੀ।

%%K t12-09-1 p
9. --- ਫਿਰ ਮੈਂ ਤੁਰਨ ਦੇ ਪਲੇਟਫ਼ਾਰਮ ਉੱਤੇ ਗਿਆ ਤੇ ਇਕ ਸਿੱਧਾ \VUgras{ਡੱਬਾ}
\textsuperscript{(37)} ਲੱਭਿਆ, ਤਾਂ ਜੋ ਰਾਹ ਵਿਚ ਬਦਲਣਾ ਨਾ ਪਵੇ। ਮੈਂ ਲਾਂਘੇ ਵਾਲੀ ਇਕ ਗੱਡੀ
ਵਿਚ ਚੜ੍ਹਿਆ, ਜਿਸ ਵਿਚ ਸਿਗਰਟ ਪੀਣ ਵਾਲਿਆਂ ਲਈ ਇਕ \VUgras{ਕੋਠੜੀ} \textsuperscript{(38)}
ਹੈ ਤੇ ਇਕ « ਇਕੱਲੀਆਂ ਬੀਬੀਆਂ » ਲਈ ਰੱਖੀ ਹੋਈ।

%%K t12-tit-2 sub
\VUpk{10.2pt}{40}{ਰਾਤ ਦਾ ਤੁਰਨਾ।}
\VUsaut{3.0mm}

%%K t12-10-1 p
10. --- \VUgras{ਪਾਏਦਾਨ} \textsuperscript{(40)} ਚੜ੍ਹ ਕੇ, \VUgras{ਬੂਹਾ}
\textsuperscript{(39)} ਬੰਦ ਕਰ ਕੇ, \VUgras{ਪਰਦਾ} \textsuperscript{(41)} ਚਾੜ੍ਹ ਕੇ ਤੇ
ਆਪਣੇ ਨਿੱਕੇ ਥੈਲੇ \VUgras{ਸਾਮਾਨ ਦੀ ਜਾਲੀ} \textsuperscript{(43)} ਵਿਚ ਰੱਖ ਕੇ ਮੈਂ
\VUgras{ਸੀਟ} \textsuperscript{(42)} ਉੱਤੇ ਬਹਿ ਗਿਆ। \VUgras{ਲੈਂਪ ਵਾਲੇ}
\textsuperscript{(44)} ਨੂੰ ਬੱਤੀਆਂ ਬਾਲਦੇ ਵੇਖ ਕੇ ਮੈਨੂੰ ਯਾਦ ਆਇਆ ਕਿ ਮੈਨੂੰ ਇਕ ਰਾਤ ਡੱਬੇ
ਵਿਚ ਲੰਘਾਉਣੀ ਹੈ, ਤੇ ਮੈਂ \VUgras{ਸਿਰਹਾਣੇ} \textsuperscript{(45)} ਕਿਰਾਏ ਉੱਤੇ ਦੇਣ ਵਾਲੇ
ਕਾਰਿੰਦੇ ਨੂੰ ਸੱਦ ਕੇ ਇਕ ਮੰਗ ਲਿਆ।

%%K t12-11-1 p
11. --- ਗੱਡੀ ਦਾ ਕੰਡਕਟਰ ਟਿਕਟ ਪਰਖਣ ਆਇਆ। ਮੈਂ ਉਸ ਕੋਲੋਂ ਪੁੱਛਿਆ ਕਿ ਅਸੀਂ ਤੁਰੇ ਕਿਉਂ
ਨਹੀਂ, ਜਦ ਕਿ ਵੇਲਾ ਹੋ ਚੁੱਕਾ ਹੈ। ਉਸ ਨੇ ਜਵਾਬ ਦਿੱਤਾ ਕਿ \VUgras{ਗੱਡੀ ਦੇ ਮੁਖੀ}
\textsuperscript{(46)} \VUgras{ਸਾਮਾਨ ਦੇ ਡੱਬੇ} \textsuperscript{(47)} ਵਿਚ ਸਾਈਕਲਾਂ
ਰਖਵਾਉਣ ਵਿਚ ਲੱਗੇ ਹਨ। ਇਸ ਤੋਂ ਬਿਨਾਂ, ਸਾਰੇ \VUgras{ਪੈਰ ਨਿੱਘੇ ਕਰਨ ਦੇ ਡੱਬੇ}
\textsuperscript{(48)} ਹਾਲੇ ਬਦਲੇ ਵੀ ਨਹੀਂ ਗਏ।

%%K t12-12-1 p
12. --- ਪਰ ਅਸੀਂ ਛੇਤੀ ਹੀ ਤੁਰਨ ਵਾਲੇ ਸੀ, ਕਿਉਂਕਿ \VUgras{ਅੱਗ ਵਾਲਾ}
\textsuperscript{(50)}, ਆਪਣੇ \VUgras{ਕੋਲੇ ਦੇ ਡੱਬੇ} \textsuperscript{(49)} ਉੱਤੇ
ਚੜ੍ਹਿਆ, \VUgras{ਇੰਜਣ} \textsuperscript{(54)} ਦੀ ਅੱਗ ਤੇਜ਼ ਕਰਨ ਨੂੰ ਕੋਲਾ ਲੈ ਰਿਹਾ ਸੀ, ਤੇ
\VUgras{ਡਰਾਈਵਰ} \textsuperscript{(51)} ਸਿਰਫ਼ \VUgras{ਨਾਇਬ ਸਟੇਸ਼ਨ ਮਾਸਟਰ}
\textsuperscript{(52)} ਦੇ ਇਸ਼ਾਰੇ ਦੀ ਉਡੀਕ ਵਿਚ ਸੀ, ਜੋ \VUgras{ਪਲੇਟਫ਼ਾਰਮ}
\textsuperscript{(53)} ਉੱਤੇ ਖੜੇ ਸਨ।

%%K t12-13-1 p
13. --- ਇੰਜਣ ਦਾ \VUgras{ਬੌਇਲਰ} \textsuperscript{(56)} ਹੌਲੀ ਸੁਰ ਵਿਚ ਗਰਜਦਾ ਸੀ;
\VUgras{ਚਿਮਨੀ} \textsuperscript{(55)} \VUgras{ਭਾਫ਼} \textsuperscript{(60)} ਰਲੇ ਧੂੰਏਂ
ਦੀਆਂ ਧਾਰਾਂ ਉਗਲਦੀ ਸੀ। ਇਕ ਤਿੱਖੀ \VUgras{ਸੀਟੀ} \textsuperscript{(59)} ਗੂੰਜੀ ਤੇ
\VUgras{ਪਿਸਟਨ} \textsuperscript{(58)} ਚੱਲਣ ਲੱਗੇ, ਜੋ ਉਸ ਭਾਰੇ ਜੰਤਰ ਦੇ \VUgras{ਪਹੀਏ}
\textsuperscript{(57)} ਘੁਮਾਉਂਦੇ ਹਨ।

%%K t12-tit-3 sub
\VUsaut{3.0mm}
\VUpk{10.2pt}{40}{ਤੁਰ ਪਏ!}
\VUsaut{3.0mm}

%%K t12-14-1 p
14. --- ਗੱਡੀ, ਜੋ \VUgras{ਪਟੜੀਆਂ} \textsuperscript{(64)} ਉੱਤੇ ਦੌੜ ਰਹੀ ਸੀ, ਤੇਜ਼ ਹੁੰਦੀ
ਗਈ; \VUgras{ਪਲੇਟਫ਼ਾਰਮ} \textsuperscript{(65)} ਮੁੱਕ ਗਏ; ਅਸੀਂ \VUgras{ਪਖ਼ਾਨਿਆਂ}
\textsuperscript{(66)} ਕੋਲੋਂ ਲੰਘੇ, ਤੇ ਇਕ ਦੂਜੀ \VUgras{ਪਟੜੀ} \textsuperscript{(63)}
ਉੱਤੇ ਖੜੇ ਡੱਬਿਆਂ ਦੀ ਕਤਾਰ ਕੋਲੋਂ ਵੀ: \VUgras{ਸੌਣ ਵਾਲੇ ਡੱਬੇ} \textsuperscript{(67)},
ਇਕ \VUgras{ਖਾਣੇ ਦਾ ਡੱਬਾ} \textsuperscript{(68)}, ਇਕ \VUgras{ਡਾਕ ਦਾ ਡੱਬਾ}
\textsuperscript{(69)}।

%%K t12-noto-1 noto
\VUnotes{143.2mm}{%
\textit{(*) ਟਿੱਪਣੀ।} --- \textit{ਇਹ ਸਮਝਿਆ ਜਾਵੇ ਕਿ ਸਫ਼ਰ ਦਾ ਬਿਆਨ ਉਹੋ ਜਿਹਾ ਹੈ ਜਿਹੋ
ਜਿਹਾ ਉਹ ਜੰਗ ਤੋਂ ਪਹਿਲਾਂ ਦੀ ਹੱਦ ਉੱਤੇ ਸੀ।}%
}

%%K t12-15-1 p
15. --- ਫਿਰ ਸਾਡੀਆਂ ਅੱਖਾਂ ਦੇ ਸਾਹਮਣੇ ਤੋਂ \VUgras{ਮਾਲ ਦਾ ਸਟੇਸ਼ਨ}
\textsuperscript{(71)} ਲੰਘਿਆ। ਛੱਪਰਾਂ ਦੇ ਹੇਠਾਂ ਕਾਰਿੰਦੇ ਹੌਲੀ ਗੱਡੀ ਨਾਲ ਘੱਲਿਆ
\VUgras{ਮਾਲ} \textsuperscript{(72)} ਲੱਦ ਰਹੇ ਸਨ। \VUgras{ਡੱਬੇ ਜੋੜਨ ਵਾਲੇ}
\textsuperscript{(76)}, \VUgras{ਪਾਣੀ ਦੀ ਟੈਂਕੀ} \textsuperscript{(73)} ਕੋਲ,
\VUgras{ਡੰਗਰਾਂ ਦੇ ਡੱਬੇ} \textsuperscript{(75)} ਤੇ \VUgras{ਖੁੱਲ੍ਹੇ ਡੱਬੇ}
\textsuperscript{(79)} ਇੱਧਰ-ਉੱਧਰ ਕਰ ਰਹੇ ਸਨ, ਇਕ \VUgras{ਮਾਲ ਗੱਡੀ}
\textsuperscript{(80)} ਬਣਾਉਣ ਨੂੰ।

%%K t12-16-1 p
16. --- ਅਸੀਂ ਆਖ਼ਰੀ \VUgras{ਕਾਂਟੇ} \textsuperscript{(78)} ਕੋਲੋਂ ਲੰਘੇ, ਜਿਸ ਕੋਲ ਕਾਂਟੇ
ਵਾਲਾ ਖੜਾ ਸੀ; ਅਸੀਂ \VUgras{ਇਸ਼ਾਰੇ ਦਾ ਚੱਕਰ} \textsuperscript{(86)} ਪਿੱਛੇ ਛੱਡਿਆ ਤੇ
ਸਟੇਸ਼ਨ ਦੂਰ ਵਿਚ ਗੁੰਮ ਹੋ ਗਿਆ।

%%K t12-17-1 p
17. --- ਛੇਤੀ ਹੀ ਅਸੀਂ ਉਸ \VUgras{ਲੋਹੇ ਦੇ ਪੁਲ} \textsuperscript{(74)} ਉੱਤੇ ਆ ਗਏ ਜੋ
\VUgras{ਦਰਿਆ} \textsuperscript{(81)} ਲੰਘਦਾ ਹੈ। ਪੁਲ ਲੰਘ ਕੇ ਅਸੀਂ ਦਰਿਆ ਦੇ ਨਾਲ ਨਾਲ ਤੁਰੇ,
ਜਿਸ ਉੱਤੇ ਤਕੜੀਆਂ \VUgras{ਖਿੱਚਣ ਵਾਲੀਆਂ ਕਿਸ਼ਤੀਆਂ} \textsuperscript{(82)} ਭਾਰੀਆਂ
\VUgras{ਮਾਲ ਕਿਸ਼ਤੀਆਂ} \textsuperscript{(83)} ਖਿੱਚ ਰਹੀਆਂ ਸਨ। ਅਸੀਂ ਇਕ
\VUgras{ਮੁਸਾਫ਼ਰ ਜਹਾਜ਼} \textsuperscript{(84)} ਵੀ ਵੇਖਿਆ, ਠੀਕ ਉਸ ਵੇਲੇ ਜਦੋਂ ਉਹ ਇਕ
\VUgras{ਤਰਦੇ ਘਾਟ} \textsuperscript{(85)} ਨਾਲ ਲੱਗਣ ਵਾਲਾ ਸੀ।

%%K t12-18-1 p
18. --- ਕਈ ਸਟੇਸ਼ਨ ਬਿਨਾਂ ਰੁਕੇ ਲੰਘ ਕੇ ਅਸੀਂ ਕੁਝ \VUgras{ਸੁਰੰਗਾਂ}
\textsuperscript{(87)} ਵਿਚੋਂ ਲੰਘੇ ਤੇ ਆਪਣੇ ਪਿੱਛੇ \VUgras{ਫਾਟਕ ਦੇ ਰਾਖਿਆਂ} ਦੀਆਂ
ਅਣਗਿਣਤ \VUgras{ਕੋਠੜੀਆਂ} \textsuperscript{(88)} ਛੱਡਦੇ ਗਏ। ਗੱਡੀ ਦੇ ਝੂਟੇ ਵਿਚ ਝੂਟੇ
ਲੈਂਦੇ ਮੈਨੂੰ ਡੂੰਘੀ ਨੀਂਦ ਆ ਗਈ।

%%K t12-tit-4 sub
\VUsaut{3.0mm}
\VUpk{10.2pt}{40}{ਡੌਇਚ-ਆਵਰੀਕੂਰ ਸਟੇਸ਼ਨ।}
\VUsaut{3.0mm}

%%K t12-19-1 p
19. --- ਮੈਨੂੰ ਅਚਾਨਕ ਇਕ ਕਾਰਿੰਦੇ ਦੀ ਆਵਾਜ਼ ਨੇ ਜਗਾ ਦਿੱਤਾ, ਜੋ ਕੂਕ ਰਿਹਾ ਸੀ:
« ਡੌਇਚ-ਆਵਰੀਕੂਰ! \textsuperscript{(*)} ਸਾਰੇ ਉਤਰੋ! » ਇਹ ਚੁੰਗੀ ਦੀ ਪੜਤਾਲ ਸੀ ਤੇ ਹੱਦ
ਦੀਆਂ ਸਾਰੀਆਂ ਝੰਜਟ ਭਰੀਆਂ ਰਸਮਾਂ। ਮੈਨੂੰ ਆਪਣੀ ਘੜੀ ਸਟੇਸ਼ਨ ਦੀ \VUgras{ਘੜੀ}
\textsuperscript{(89)} ਨਾਲ ਮਿਲਾਉਣੀ ਪਈ, ਜੋ ਪੰਜਾਹ ਪੰਜ ਮਿੰਟ ਅੱਗੇ ਸੀ; ਮੁਸ਼ਕਲ ਨਾਲ ਜਾਗਿਆ
ਹੋਇਆ ਮੈਂ \VUgras{ਵੱਡੀ ਸੂਈ} \textsuperscript{(90)} ਤੇ \VUgras{ਨਿੱਕੀ}
\textsuperscript{(91)} ਵਿਚ ਫ਼ਰਕ ਔਖਾ ਹੀ ਕਰ ਸਕਦਾ ਸੀ; \VUgras{ਡਾਇਲ}
\textsuperscript{(92)} ਤਾਂ ਸਵੇਰ ਦੀ ਧੁੰਦ ਨਾਲ ਬਿਲਕੁਲ ਧੁੰਦਲਾ ਸੀ।

%%K t12-20-1 p
20. --- ਜਦੋਂ ਤਕ ਚੁੰਗੀ ਵਾਲੇ ਆਪਣੀ ਪੜਤਾਲ ਕਰਦੇ ਰਹੇ, ਮੈਂ ਉਬਾਸੀਆਂ ਲੈਂਦੇ ਹੋਏ ਉਨ੍ਹਾਂ
\VUgras{ਇਸ਼ਤਿਹਾਰਾਂ} \textsuperscript{(93)} ਨੂੰ ਪੜ੍ਹਦਾ ਰਿਹਾ ਜੋ ਸਟੇਸ਼ਨ ਦੀਆਂ ਕੰਧਾਂ
ਸਜਾਉਂਦੇ ਹਨ। ਵੇਲਾ ਕੱਟਣ ਨੂੰ ਮੈਂ ਨਾਸ਼ਤਾ ਕੀਤਾ, ਜਰਮਨ \VUgras{ਜਾਲ} \textsuperscript{(94)}
ਦਾ ਨਕਸ਼ਾ ਵੇਖਿਆ ਕਿ ਸਟ੍ਰਾਸਬੂਰਗ ਹਾਲੇ ਕਿੰਨੀ ਦੂਰ ਹੈ, ਤੇ ਆਪਣੇ ਸਫ਼ਰ ਦੇ ਅਖ਼ੀਰ ਤਕ ਛੇਤੀ
ਪਹੁੰਚਣ ਦੀ ਆਸ ਨਾਲ ਦਿਲ ਵਿਚ ਜ਼ੋਰ ਪਾ ਕੇ ਫਿਰ ਡੱਬੇ ਵਿਚ ਚੜ੍ਹ ਗਿਆ।
\VUsaut{6.0mm}
\VUfilet{22mm}
